% !TEX root = main.tex
%% Figure & Graphics %%%%%%%%%%%%%%%%%%%%%%%%%%%%%%%%%%%%%%%%%%%%%%%%%%%%%%%%%%
\usepackage{graphicx}
\usepackage[dvipsnames]{xcolor}
{% raw -%}
%\graphicspath{{subdir1/}{../dir2}}
{% endraw -%}
%
% Nord colors
\definecolor{night1}{rgb}{0.180, 0.204, 0.251}
\definecolor{night2}{rgb}{0.231, 0.259, 0.322}
\definecolor{night3}{rgb}{0.263, 0.298, 0.369}
\definecolor{night4}{rgb}{0.298, 0.337, 0.416}
\definecolor{snow1}{rgb}{0.847, 0.871, 0.914}
\definecolor{snow2}{rgb}{0.898, 0.914, 0.941}
\definecolor{snow3}{rgb}{0.925, 0.937, 0.957}
\definecolor{frost1}{rgb}{0.561, 0.737, 0.733}
\definecolor{frost2}{rgb}{0.533, 0.753, 0.816}
\definecolor{frost3}{rgb}{0.506, 0.631, 0.757}
\definecolor{frost4}{rgb}{0.369, 0.506, 0.675}
\definecolor{nordred}{rgb}{0.749, 0.380, 0.416}
\definecolor{nordorange}{rgb}{0.816, 0.529, 0.439}
\definecolor{nordyellow}{rgb}{0.922, 0.796, 0.545}
\definecolor{nordgreen}{rgb}{0.639, 0.745, 0.549}
\definecolor{nordviolet}{rgb}{0.706, 0.557, 0.678}
%
%%%% Math expression %%%%%%%%%%%%%%%%%%%%%%%%%%%%%%%%%%%%%%%%%%%%%%%%%%%%%%%%%%
{% raw -%}
\usepackage{amsmath,amssymb}
\usepackage{mathtools}
\usepackage{bm}
\usepackage{braket}
\usepackage{siunitx}
\sisetup{input-digits = 0123456789\pi}
% \mathtoolsset{showonlyrefs=true} % Show equation labels only referenced
\DeclareMathOperator{\Tr}{Tr}
\renewcommand{\Re}{\operatorname{Re}}
\renewcommand{\Im}{\operatorname{Im}}
\newcommand{\defas}{\equiv}
% theorem like environment
\usepackage{amsthm}
% --- plainbreak: plain の改行版（本文イタリック，見出し太字） ---
\newtheoremstyle{plainbreak}%
  {\topsep}%   見出し上のスペース
  {\topsep}%   本文下のスペース
  {\itshape}%  本文フォント
  {}%          インデント
  {\bfseries}% 見出しフォント
  {}%         見出し末尾
  {\newline}%  見出し後に改行
  {\thmname{#1}\thmnumber{ #2}\thmnote{: #3}}%
% --- definitionbreak: definition の改行版（本文立体，見出し太字） ---
\newtheoremstyle{definitionbreak}%
  {\topsep}%
  {\topsep}%
  {\upshape}%  本文は立体
  {}%
  {\bfseries}%
  {}%
  {\newline}%
  {\thmname{#1}\thmnumber{ #2}\thmnote{: #3}}%
% --- remarkbreak: remark の改行版（本文立体，見出しイタリック，スペース控えめ） ---
\newtheoremstyle{remarkbreak}%
  {0.5\topsep}%  remark 系は上下スペースが半分
  {0.5\topsep}%
  {\upshape}%
  {}%
  {\itshape}%    見出しはイタリック（太字でない）
  {}%
  {\newline}%
  {\thmname{#1}\thmnumber{ #2}\thmnote{: #3}}%
% --- lawbreak: 法則用（タイトルのみを見出しに，番号なし，見出し後改行） ---
\newtheoremstyle{lawbreak}%
  {\topsep}%
  {\topsep}%
  {\upshape}%   本文は立体
  {}%
  {\bfseries}%  見出し太字
  {}%           句読点なし（タイトルが見出しそのものなので）
  {\newline}%
  {\thmnote{#3}}%   ← note だけを見出しとして表示
% --- 定理本文がリスト環境で始まっても見出し後の改行を効かせる ---
%     amsthm の \newline スタイルは見出しを箱 \@labels に貯めて次の段落開始時の
%     \everypar（\dth@everypar）で取り付けるが，本文が直接リストで始まると
%     カーネルの \@donoparitem が見出し箱を最初の \item と同じ行に連結して
%     しまい，箱の中に封じられた改行ペナルティ（\nobreak\hfil\break）は効かない．
%     対策は2段構え:
%     (1) \thmheadnl と \hskip\thm@headsep を見出し箱の中に入れず，段落開始時に
%         箱の外で実行する（\@begintheorem と \dth@everypar を再定義）．
%         通常のテキスト本文では出力は従来と同一．
%     (2) リスト環境の開始直前に見出しが保留中なら，\thmheadnl を一時的に
%         無効化した上で \leavevmode で段落を発生させて見出しを取り付ける
%         （改行ノードは最初から出さない．LuaTeX-ja が箱や \hskip の周囲に
%         whatsit を挿入するため，\unskip\unpenalty による事後除去は不可能）．
%         直後にリストが段落を閉じるので見出しが1行を成し，リストはその下から
%         始まる．
\makeatletter
\newif\ifcallout@thmhead
\let\callout@orig@deferredthmhead\deferred@thm@head
\def\deferred@thm@head#1{%
  \global\callout@thmheadtrue
  \callout@orig@deferredthmhead{#1}%
}
% amsthm の \@begintheorem から \thmheadnl \hskip\thm@headsep を箱の外へ移設
\def\@begintheorem#1#2[#3]{%
  \deferred@thm@head{\the\thm@headfont \thm@indent
    \@ifempty{#1}{\let\thmname\@gobble}{\let\thmname\@iden}%
    \@ifempty{#2}{\let\thmnumber\@gobble}{\let\thmnumber\@iden}%
    \@ifempty{#3}{\let\thmnote\@gobble}{\let\thmnote\@iden}%
    \thm@swap\swappedhead\thmhead{#1}{#2}{#3}%
    \the\thm@headpunct
  }%
  \ignorespaces}
% amsthm の \dth@everypar に移設分の実行とフラグ解除を追加
\dth@everypar={%
  \@minipagefalse \global\@newlistfalse
  \@noparitemfalse
  \if@inlabel
    \global\@inlabelfalse
    \begingroup \setbox\z@\lastbox
     \ifvoid\z@ \kern-\itemindent \fi
    \endgroup
    \unhbox\@labels
    \thmheadnl\hskip\thm@headsep % ← 箱の外で実行（移設分）
  \fi
  \if@nobreak \@nobreakfalse \clubpenalty\@M
  \else \clubpenalty\@clubpenalty \everypar{}%
  \fi
  \global\callout@thmheadfalse % ← 見出しが取り付いたらフラグを下ろす
}%
\def\callout@nlmarker{\newline}
\newcommand*{\callout@fixthmheadbreak}{%
  \ifcallout@thmhead
    \ifx\thmheadnl\callout@nlmarker % 改行スタイルの見出しのみ対象
      \begingroup
        \let\thmheadnl\relax
        \thm@headsep\z@skip
        \leavevmode % 改行抜きで見出しを取り付ける（段落はリストが閉じる）
      \endgroup
    \fi
  \fi
}
\AddToHook{env/itemize/before}{\callout@fixthmheadbreak}
\AddToHook{env/enumerate/before}{\callout@fixthmheadbreak}
\AddToHook{env/description/before}{\callout@fixthmheadbreak}
% verbatim（カーネル）も内部はリスト（trivlist）なので機構は同一．
% Verbatim は fancyvrb / fvextra / minted を使う場合への備え
% （これらを読み込んでいなくても，未定義環境への \AddToHook は無害）
\AddToHook{env/verbatim/before}{\callout@fixthmheadbreak}
\AddToHook{env/Verbatim/before}{\callout@fixthmheadbreak}
\makeatother
% Law env
\theoremstyle{lawbreak}
\newtheorem*{law@inner}{Law}% 環境名は内部用．"Law" は表示されない（\thmname を使わないため）
\makeatletter
\NewDocumentEnvironment{law}{ m }{%
  \begin{law@inner}[#1]%
}{%
  \end{law@inner}%
}
\makeatother
% Envs with number
\theoremstyle{definitionbreak}
%\theoremstyle{definition} % 改行なしスタイル
\newtheorem{thm}{定理}[section]
\newtheorem{lem}[thm]{補題}
\newtheorem{prop}[thm]{命題}
\newtheorem{cor}[thm]{系}
\newtheorem{conj}[thm]{予想}
\newtheorem{asm}[thm]{仮定}
\newtheorem{dfn}[thm]{定義}
\newtheorem{rem}[thm]{注}
\newtheorem{exc}{練習問題}[section]
%
% Envs without number
\newtheorem*{thm*}{定理}
\newtheorem*{lem*}{補題}
\newtheorem*{prop*}{命題}
\newtheorem*{cor*}{系}
\newtheorem*{conj*}{予想}
\newtheorem*{asm*}{仮定}
\newtheorem*{dfn*}{定義}
\newtheorem*{rem*}{注}
\newtheorem*{exc*}{練習問題}
{% endraw -%}
%%%% Callout block %%%%%%%%%%%%%%%%%%%%%%%%%%%%%%%%%%%%%%%%%%%%%%%%%%%%%%%%%%%%
{% raw -%}
\usepackage[most]{tcolorbox}
% --- ボックス内ではリストの余白を詰める ---
%     クラスの既定（ltjsarticle: \leftmargini=3\zw, \topsep=0.5\baselineskip，
%     scrartcl 等: \leftmargini=2.5em + 大きめの topsep/itemsep）は
%     パディングの狭いボックス内では過大なため，第1レベルのみ局所的に上書きする
\makeatletter
\newcommand*{\callout@tightlists}{%
  \ifdefined\zw \leftmargini=1.5\zw \else \leftmargini=1.5em \fi % \zw は LuaTeX-ja 専用単位
  \def\@listi{\leftmargin\leftmargini
    \labelwidth\leftmargini \advance\labelwidth-\labelsep
    \parsep\z@
    \topsep 0.2\baselineskip
    \itemsep\z@\relax}%
}
\tcbset{
  calloutbase/.style={
    arc=0.7mm,
    boxrule=0mm,
    left=2.5mm,
    leftrule=0.7mm,
    right=3.5mm,
    top=1.5mm,
    bottom=1.5mm,
    colback=snow3,
    before upper={\callout@tightlists},
  },
  callout/.style 2 args={   % #1: タイトル, #2: タイトル下スペース
    calloutbase,
    coltitle=black,
    fonttitle=\bfseries,
    detach title,
    before upper={\callout@tightlists\tcbtitle\par\vspace{#2}},
    title={#1},
  },
}
\makeatother
% --- calloutbox: タイトル付きの囲み領域（直接使う用） ---
\NewTColorBox{calloutbox}{ O{night4} o O{1mm} }{%
  colframe=#1,
  IfValueTF={#2}{callout={#2}{#3}}{calloutbase}%
}
% --- 対応表: 環境名 → 枠色 ---
\makeatletter
\NewDocumentCommand{\RegisterCalloutEnv}{ m m O{optional} }{%
  % #1: 環境名, #2: 色, #3: タイトル渡し方式 optional / mandatory
  \@namedef{callout@color@#1}{#2}%
  \@namedef{callout@titlemode@#1}{#3}%
}
\makeatother
% 環境と色の対応関係
\RegisterCalloutEnv{thm}{frost1}
\RegisterCalloutEnv{lem}{frost1}
\RegisterCalloutEnv{prop}{night4}
\RegisterCalloutEnv{cor}{night4}
\RegisterCalloutEnv{conj}{night4}
\RegisterCalloutEnv{thm*}{frost1}
\RegisterCalloutEnv{lem*}{frost1}
\RegisterCalloutEnv{prop*}{night4}
\RegisterCalloutEnv{cor*}{night4}
\RegisterCalloutEnv{conj*}{night4}
%
\RegisterCalloutEnv{dfn}{nordviolet}
\RegisterCalloutEnv{asm}{nordviolet}
\RegisterCalloutEnv{dfn*}{nordviolet}
\RegisterCalloutEnv{asm*}{nordviolet}
%
\RegisterCalloutEnv{law}{nordviolet}[mandatory]
%
\RegisterCalloutEnv{rem}{night4}
\RegisterCalloutEnv{rem*}{night4}
%
\RegisterCalloutEnv{exc}{nordgreen}
\RegisterCalloutEnv{exc*}{nordgreen}
%
% --- callout: 定理系環境を囲む用 ---
%     省略可能タイトルは内部環境のオプショナル引数（定理名）として渡す
\makeatletter
\NewDocumentEnvironment{callout}{ o m }{%
  \ifcsname callout@color@#2\endcsname\else
    \GenericError{}{callout環境: `#2' は未登録です．
      \RegisterCalloutEnv で登録してください}{}{}%
  \fi
  \begin{calloutbox}[\@nameuse{callout@color@#2}]%
  \ifcsstring{callout@titlemode@#2}{mandatory}
    {% タイトル必須の環境
     \IfValueTF{#1}
       {\begin{#2}{#1}}
       {\GenericError{}{callout環境: `#2' はタイトル必須です}{}{}%
        \begin{#2}{??}}%
    }
    {% タイトル省略可能の環境
     \IfValueTF{#1}
       {\begin{#2}[#1]}
       {\begin{#2}}%
    }%
}{%
  \end{#2}%
  \end{calloutbox}%
}
\makeatother
%
\NewDocumentCommand{\NewCalloutBoxEnv}{ m m o }{%
  % #1: 環境名, #2: 色, #3: デフォルトタイトル（省略可）
  \IfValueTF{#3}
    {% デフォルトタイトルあり: 環境のタイトル省略時に #3 を使う
     \NewDocumentEnvironment{#1}{ O{#3} }{%
       \begin{calloutbox}[#2][##1]%
     }{%
       \end{calloutbox}%
     }%
    }
    {% デフォルトタイトルなし: 省略時はタイトルなしのボックス
     \NewDocumentEnvironment{#1}{ o }{%
       \IfValueTF{##1}
         {\begin{calloutbox}[#2][##1]}
         {\begin{calloutbox}[#2]}%
     }{%
       \end{calloutbox}%
     }%
    }%
}
%
\NewCalloutBoxEnv{objective}{nordred}[目標]
\NewCalloutBoxEnv{summary}{nordred}[まとめ]
\NewCalloutBoxEnv{info}{nordyellow}
{% endraw -%}
%
%% Bibliography %%%%%%%%%%%%%%%%%%%%%%%%%%%%%%%%%%%%%%%%%%%%%%%%%%%%%%%%%%%%%%%
{% if cookiecutter.bibliography == "biblatex+bibtex" -%}
\usepackage[backend=bibtex, style=phys, biblabel=brackets,
           language=auto, autolang=langname,
           articletitle=false, pageranges=false,
           maxnames=3, maxcitenames=2]{biblatex}
\DefineBibliographyStrings{english}{
   andothers = {\mkbibemph{et\addabbrvspace al\adddot}} % Print et al. in italic
}
%\addbibresource{<bibfile1>.bib} % The extension .bib is required!
%\addbibresource{<bibfile2>.bib}
{% elif cookiecutter.bibliography == "biblatex+biber" -%}
\usepackage[backend=biber, style=phys, biblabel=brackets,
           language=auto, autolang=langname,
           articletitle=false, pageranges=false,
           maxnames=3, maxcitenames=2]{biblatex}
\DefineBibliographyStrings{english}{
   andothers = {\mkbibemph{et\addabbrvspace al\adddot}} % Print et al. in italic
}
% \addbibresource{<bibfile1>.bib} % The extension .bib is required!
% \addbibresource{<bibfile2>.bib}
{% elif cookiecutter.bibliography == "bibtex" -%}
\bibliographystyle{jplain} 
% stylename=jabbrv,jalpha,jplain, or junsrt
{% endif -%}
{% if cookiecutter.bibliography == "bibtex" or cookiecutter.bibliography == "manual" -%}
\usepackage{cite}
{% endif -%}
%
%% Settings for margins %%%%%%%%%%%%%%%%%%%%%%%%%%%%%%%%%%%%%%%%%%%%%%%%%%%%%%%
%\usepackage{geometry}
%\geometry{left=30truemm,right=30truemm,textheight=25truecm}
%
%% Misc %%%%%%%%%%%%%%%%%%%%%%%%%%%%%%%%%%%%%%%%%%%%%%%%%%%%%%%%%%%%%%%%%%%%%%%
{% raw -%}
\usepackage{booktabs}
\usepackage[hang,flushmargin,perpage]{footmisc}
% hang=インデントを除去，flushmargin=記号と脚注のスペースを除去，perpage=ページごとに独立した脚注番号
\setlength{\footnotemargin}{1em} % 脚注記号と脚注の間のスペースを調整
\usepackage{algorithm,algpseudocode}
\usepackage{url}
% \usepackage[%
%     %dvipdfmx,
%     %luatex,
%     hidelinks,%
%     colorlinks=true,%
%     urlcolor=frost4,%
%     linkcolor=frost4,%
%     citecolor=nordgreen,%
% ]{hyperref}
{% endraw -%}
%
%% Draft %%%%%%%%%%%%%%%%%%%%%%%%%%%%%%%%%%%%%%%%%%%%%%%%%%%%%%%%%%%%%%%%%%%%%%
{% raw -%}
%\usepackage[color]{showkeys}
%\renewcommand*{\showkeyslabelformat}[1]{%
%\fbox{\parbox{40truemm}{\normalfont\small\ttfamily#1}}}
%\usepackage[]{setspace}
%\doublespacing
%\newcommand*{\TODO}[1]{{\color{red}(TODO: #1)}}
%\newcommand{\remove}[1]{\textcolor{green}{#1}}
%\newcommand{\add}[1]{\textcolor{Purple}{#1}}
% Replace figures with dummy figures
%\let\originalincludegraphics\includegraphics
%\renewcommand{\includegraphics}[2][width=3.2truein]{\originalincludegraphics[#1]{example-image-plain.pdf}}
{% endraw -%}

